CanSat on Euroopa Kosmoseagentuuri poolt korraldatud võistlus, mis on suunatud noortele kooliõpilastele. Põhimõtteliselt seisneb võistluse eesmärk selles, et luua satelliit, mis oleks limonaadipurgisuurune ning saaks edastada andmeid anduritelt maal oleva jaamale. Osalejad on vabad valida, kuidas nad konstrueerivad seda. Meeskond, mis võidab, saab auhinnaks reisi koos juhendajaga Hollandisse ESA ESTEC kosmosekeskusesse. Sellel peab olema töötav langevari, et ennetada satelliidi purunemist, ning ka vastupidav korpus, mis hoiaks elektroonikat sees hästi ning annaks ligipääsu õhule, samal ajal ennetades purunemist. Osalejad peavad ise seda kokku panema, programmeerima ja dekoreerima.

CanSati võistluste meeskonnas võib olla 4–6 õpilast, keda abistab õpetaja või juhendaja. Õpilaste vanus on 14-19 aastat.

CanSati võistlus toimub Eestis 24.04.2025, kus meeskondade poolt ehitatud satelliidid heidetakse alla Tallinna Teletornist. Langevarjuga alla laskumise ajal peab satelliit täitma kaks missiooni:
\begin{enumerate}
    \item CanSat peab mõõtma õhutemperatuuri ja õhurõhku reaalajas ning edastada andmed maapealsele jaamale (mõõdetud andmed tuleb edastada iga sekund).
    \item Teise missiooni eesmärgi valivad meeskonnad ise. Meie meeskond valis  
    õhuniiskuse ja CO$_2$ sisalduse õhus mõõtmise SCD41 anduri abil.
\end{enumerate}

